% Examples shared by the Y statements and their individual proofs.
\newcommand{\Yexampleone}{%
\textbf{Y1: one or two switches}\par\medskip
\begin{shiftchain}[
  column sep=5mm, row sep=6mm, dot radius=3pt,
  labels=false, label font=\scriptsize,
  bars={3,7}, bar labels={3/$y_j$,7/$y_{j+1}$},
]{8}
  \shiftvertexcolor{1,2,8}{cfred}
  \shiftvertexcolor{4,5,6}{cfnavy}
  \shiftedge{1,2,4,5}[$\vec e$]
  \shiftedge{2,4,6,8}[$\vec g$]
\end{shiftchain}
\par\medskip
$\vec e$: one switch.\quad $\vec g$: two switches.
}

\newcommand{\Yexampletwo}{%
\textbf{Y2: each switch stays in its interval}\par\medskip
\begin{shiftchain}[
  column sep=5mm, row sep=6mm, dot radius=3pt,
  labels=false, label font=\scriptsize,
  bars={2/dashed,5,8/dashed},
  bar labels={2/$x_{j-1}$,5/$y_j(t)$,8/$x_j$},
]{9}
  \shiftedge{3/cfred,5/cfred,7/cfnavy,9/cfnavy}[$\vec e$]
\end{shiftchain}
\par\medskip
Dashed bars: original positions.
}

\newcommand{\Yexamplethree}{%
\textbf{Y3: the next switch is past $e_2$}\par\medskip
\begin{shiftchain}[
  column sep=5mm, row sep=6mm, dot radius=3pt,
  labels=false, label font=\scriptsize,
  bars={3,5/dashed,7},
  bar labels={3/$y_j(t)$,5/$x_j$,7/$y_{j+1}(t)$},
]{9}
  \foreach \pos/\lab in {1/e_1,4/e_2,6/e_3,9/e_4}
    \node at (\pos,0) {$\lab$};
  \shiftedge{1/cfred,4/cfnavy,6/cfnavy,9/cfred}[$\vec e$]
\end{shiftchain}
\par\medskip
$e_1<x_j\quad\Longrightarrow\quad e_2<y_{j+1}(t)$.
}

\newcommand{\Yexamplefour}{%
\textbf{Y4: forbidden pattern}\par\medskip
\begin{shiftchain}[
  column sep=5mm, row sep=6mm, dot radius=3pt,
  labels=false, label font=\scriptsize,
  bars={2,4}, bar labels={2/$y_j$,4/$y_{j+1}$},
]{9}
  \foreach \pos/\lab in {1/e_1,5/e_2,7/e_3,9/e_4}
    \node at (\pos,0) {$\lab$};
  \shiftedge[cfred]{1,5,7,9}[$\vec e$]
\end{shiftchain}
\par\medskip
\textcolor{cfred}{No two switches in $(e_1,e_2)$.}
}

\newcommand{\Yexamplefive}{%
\textbf{Y5: a witness for a $(2,3)$-edge}\par\medskip
\begin{shiftchain}[
  column sep=4mm, row sep=7mm, dot radius=3pt,
  labels=false, label font=\scriptsize,
  bars={3,6,8/dashed,10},
  bar labels={3/$y_{j-1}$,6/$y_j$,8/$x_j$,10/$y_{j+1}$},
]{12}
  \node at (7,0) {$e_3=g_3$};
  \shiftvertexcolor{1,2,7,9}{cfred}
  \shiftvertexcolor{4,5,12}{cfnavy}
  \shiftedge{1,2,7,9}[$\vec e$]
  \shiftedge{4,5,7,12}[$\vec g$]
\end{shiftchain}
\par\medskip
$e_3^-=y_j=g_3^-<x_j$\par\smallskip
$y_{j+1}<g_4$.
}

\begin{frame}[t]{Stage Y: recap}
  \small\setstretch{1.0}
  \claimreminder{\Xthree}
  Sweep right to left; move the first switch of each $(3,4)$-edge
  just before its third vertex.

  \smallskip
  \centering
  % Key steps from the algorithm outline, on the same numbered chain.
  \only<1>{%
    \outlinediagram
      {3,6,9,13,16,19}
      {3/$x_1$,6/$x_2$,9/$x_3$,13/$x_4$,16/$x_5$,19/$x_6$}
      {1,2,3,7,8,9,14,15,16,20}
      {4,5,6,10,11,12,13,17,18,19}
      {9}
      {Y: right to left (RL) --- $e_{9}$ is monochromatic}
      {-1}%
  }%
  \only<2>{%
    \outlinediagram
      {3,6,9,13,15/outline moved,16/outline old,19}
      {3/$x_1$,6/$x_2$,9/$x_3$,13/$x_4$,15/$y_5$,19/$x_6$,16/$x_5$}
      {1,2,3,7,8,9,14,15,20}
      {4,5,6,10,11,12,13,16,17,18,19}
      {9}
      {Y: right to left (RL) --- move $y_5$ before vertex $16$}
      {-1}%
  }%
  \only<3>{%
    \outlinediagram
      {3,6,9,13,15,19}
      {3/$x_1$,6/$x_2$,9/$x_3$,13/$x_4$,15/$y_5$,19/$x_6$}
      {1,2,3,7,8,9,14,15,20}
      {4,5,6,10,11,12,13,16,17,18,19}
      {4}
      {Y: right to left (RL) --- $e_{4}$ is monochromatic}
      {-1}%
  }%
  \only<4>{%
    \outlinediagram
      {3,5/outline moved,6/outline old,9,13,15,19}
      {3/$x_1$,5/$y_2$,9/$x_3$,13/$x_4$,15/$y_5$,19/$x_6$,6/$x_2$}
      {1,2,3,6,7,8,9,14,15,20}
      {4,5,10,11,12,13,16,17,18,19}
      {4}
      {Y: right to left (RL) --- move $y_2$ before vertex $6$}
      {-1}%
  }%
  \only<5>{%
    \outlinediagram
      {3,5,9,13,15,19}
      {3/$x_1$,5/$y_2$,9/$x_3$,13/$x_4$,15/$y_5$,19/$x_6$}
      {1,2,3,6,7,8,9,14,15,20}
      {4,5,10,11,12,13,16,17,18,19}
      {1}
      {Y: right to left (RL) --- $e_{1}$}
      {-1}%
  }%
\end{frame}

% Introduce every claim before starting the induction.  The statements stay
% visible as the example on the right changes to the newly introduced claim.
\begin{frame}[t]{Y1--Y6: what Stage Y preserves and achieves}
  {\normalfont\small\color{black}\setstretch{1}
    $t$: number of Stage Y switch moves performed\par
    $c_{Y,t}$: colouring after $t$ moves\par
    $c_{Y,0}=c_X$, \quad $c_{Y,T}=c_Y$ \quad ($T$: total number of moves)\par
    Induction on $t$: Y1--Y5 preserved; Y6 at termination\par
  }
  \smallskip
  \begin{columns}[T,onlytextwidth]
    \begin{column}{0.58\textwidth}
      \footnotesize\setstretch{1}
      \begin{itemize}[itemsep=5pt]
        \item<1-> \Yone
        \item<2-> \Ytwo
        \item<3-> \Ythree
        \item<4-> \Yfour
        \item<5-> \Yfive
        \item<6-> \Ysix
      \end{itemize}
    \end{column}
    \begin{column}{0.39\textwidth}
      \centering\small\setstretch{1}
      \begin{overlayarea}{\linewidth}{5.5cm}
        \centering
        \only<1>{\Yexampleone}
        \only<2>{\Yexampletwo}
        \only<3>{\Yexamplethree}
        \only<4>{\Yexamplefour}
        \only<5>{\Yexamplefive}
        \only<6>{
          \textbf{Y6: the final monochromatic pattern}\par\medskip
          \begin{shiftchain}[
            column sep=5mm, row sep=6mm, dot radius=3pt,
            labels=false, label font=\scriptsize,
            bars={3,6}, bar labels={3/$y_{j-1}$,6/$y_j$},
          ]{9}
            \foreach \pos/\lab in {1/e_1,2/e_2,7/e_3,9/e_4}
              \node at (\pos,0) {$\lab$};
            \shiftedge[cfred]{1,2,7,9}[$\vec e$]
          \end{shiftchain}
          \par\medskip
          In $c_Y$: both switches in $(e_2,e_3)$.
        }
      \end{overlayarea}
    \end{column}
  \end{columns}
\end{frame}

% Inductive step of Lemma 3.2 in Smith.
\begin{frame}{Y1--Y5: one Stage Y step}
  Induction on $t$: assume Y1--Y5 at time $t$.\par
  Choose a maximal $(3,4)$-edge $\vec f$.
  Move only $y_i(t)$, setting $y_i(t+1)=f_3^-$.\par

  \begin{center}
    \only<1>{
      \begin{shiftchain}[
        column sep=6mm, row sep=5mm, dot radius=3.0pt,
        label font=\scriptsize,
        bars={0,8,10,12},
        bar labels={0/$y_{i-1}(t)$,8/\raisebox{-3.5mm}{$y_i(t)$},10/$y_{i+1}(t)$,12/$y_{i+2}(t)$},
      ]{12}
        \shiftedge{2/cfred,4/cfred,7/cfred,12/cfred}[$\vec f$]
      \end{shiftchain}
    }
    \only<2->{
      \begin{shiftchain}[
        column sep=6mm, row sep=5mm, dot radius=3.0pt,
        label font=\scriptsize,
        bars={0,6,8/dashed,10,12},
        bar labels={0/$y_{i-1}(t)$,6/$f_3^-$,8/\raisebox{-3.5mm}{$y_i(t)$},10/$y_{i+1}(t)$,12/$y_{i+2}(t)$},
      ]{12}
        \shiftedge{2/cfred,4/cfred,7/cfnavy,12/cfred}[$\vec f$]
      \end{shiftchain}
    }
  \end{center}

  \begin{itemize}
    \item<1-> $y_{i-1}(t)<f_1<f_3<y_i(t)<y_{i+1}(t)<f_4<y_{i+2}(t)$.
    \item<2-> $y_{i-1}(t)<f_1<f_2<y_i(t+1)=f_3^-<f_3<y_{i+1}(t)<f_4$.
  \end{itemize}
\end{frame}

\begin{frame}[t]{Auxillary claim: the moved edge starts after $x_{i-1}$}
  \claimreminder{\Ytwo\Xtwo}

  \begin{lemma}{Claim}
    $x_{i-1}<f_1$
  \end{lemma}

  \begin{proof}{Proof}
    \begin{itemize}
      \item<1-> Suppose $f_1<x_{i-1}$.
      \item<2-> Y2 gives $x_i<y_{i+1}(t)<f_4$,
        so $\vec f$ contains $x_{i-1},x_i$.
      \item<3-> X2 forces $x_i=f_4^-$, contradicting
        $x_i<y_{i+1}(t)<f_4$.
    \end{itemize}
  \end{proof}

  \begin{center}
    \only<-2>{
      \begin{shiftchain}[
        column sep=6mm, row sep=5mm, dot radius=3.0pt,
        labels=false, label font=\scriptsize,
        bars={3/dashed,7,9/dashed,11/dashed,13},
        bar labels={3/$x_{i-1}$,7/$y_i(t+1)$,9/\raisebox{-3.5mm}{$y_i(t)$},11/$x_i$,13/$y_{i+1}(t)$},
      ]{16}
        \shiftedge{2/cfred,5/cfred,8/cfnavy,16/cfred}[$\vec f$]
      \end{shiftchain}
    }
    \only<3->{
      \begin{shiftchain}[
        column sep=6mm, row sep=5mm, dot radius=3.0pt,
        labels=false, label font=\scriptsize,
        bars={3/dashed,7,9/dashed,11/dashed,13,15/densely dotted},
        bar labels={3/$x_{i-1}$,7/$y_i(t+1)$,9/\raisebox{-3.5mm}{$y_i(t)$},11/$x_i$,13/$y_{i+1}(t)$,15/$f_4^-$},
      ]{16}
        \shiftedge{2/cfred,5/cfred,8/cfnavy,16/cfred}[$\vec f$]
      \end{shiftchain}
    }
  \end{center}
\end{frame}


% Prove each invariant from its base case and inductive step.
% Y5 is split into separate frames for the possible switch pairs.
\begin{frame}[t]{Y1: every edge keeps one or two switches}
  \footnotesize\setstretch{1}
  \tcbset{before skip=0pt,after skip=0pt}
  \claimreminder{\Xone\Xtwo\Ytwo
    \textbf{Auxiliary claim.} $x_{i-1}<f_1$.}
  \begin{lemma}{Claim}\Yone\end{lemma}
  \begin{proof}{Proof}
    \begin{itemize}[itemsep=0pt,topsep=0pt,parsep=0pt]
      \item<1-> \textbf{Base.} $c_{Y,0}=c_X$: X1.
      \item<2-> \textbf{Step: no zero.} Losing the only bar gives
        $y_i(t+1)<e_1<y_i(t)<e_4<y_{i+1}(t)$.
        Thus $f_1<e_1$ but $e_4<f_4$: incomparable.
        \only<2->{%
          \par\vspace{-1pt}
          {\centering
            \begin{shiftchain}[
              column sep=5.5mm, row sep=3.5mm, dot radius=2.7pt,
              labels=false, label font=\scriptsize,
              tikz options={scale=0.9,transform shape},
              bars={6,9/dashed,11},
              bar labels={6/$y_i(t+1)$,9/\raisebox{-3.5mm}{$y_i(t)$},11/$y_{i+1}(t)$},
            ]{12}
              \shiftedge{2/cfred,4/cfred,7/cfnavy,12/cfred}[$\vec f$]
              \shiftedge{7/cfnavy,8/cfnavy,9/cfnavy,10/cfnavy}[$\vec e$]
            \end{shiftchain}
          \par}
        }
      \item<3-> \textbf{No third bar.} A new third bar $y_i(t+1)=f_3^-$
        joins $y_{i-2}(t),y_{i-1}(t)$. By Y2,
        $e_1<y_{i-2}(t)\leq x_{i-2}<y_{i-1}(t)\leq x_{i-1}$.
        \only<3->{%
          \par\vspace{-1pt}
          {\centering
            \begin{shiftchain}[
              column sep=5.5mm, row sep=3.5mm, dot radius=2.7pt,
              labels=false, label font=\scriptsize,
              tikz options={scale=0.9,transform shape},
              bars={2,5,9,12/dashed,13},
              bar labels={2/$y_{i-2}(t)$,5/$y_{i-1}(t)$,9/$y_i(t+1)$,12/\raisebox{-3.5mm}{$y_i(t)$}},
            ]{14}
              \shiftedge{7/cfred,8/cfred,10/cfnavy,14/cfred}[$\vec f$]
              \shiftedge{1/cfred,4/cfnavy,7/cfred,10/cfnavy}[$\vec e$]
            \end{shiftchain}
          \par}
        }
      \item<4-> Since $x_{i-1}<f_1<f_3\leq e_4$, X2 gives
        $x_{i-1}=e_4^->f_1$: a contradiction.
    \end{itemize}
  \end{proof}
\end{frame}

\begin{frame}[t]{Y2: each switch stays in its interval}
  \footnotesize\setstretch{1}
  \tcbset{before skip=4pt,after skip=4pt}
  \claimreminder{\textbf{Auxiliary claim.} $x_{i-1}<f_1$.}
  \begin{lemma}{Claim}\Ytwo\end{lemma}
  \begin{columns}[T,onlytextwidth]
    \begin{column}{0.61\textwidth}
      \begin{proof}{Proof}
        \begin{itemize}[itemsep=5pt]
          \item<1-> \textbf{Base.} $y_j(0)=x_j$.
          \item<2-> \textbf{Step.} Only $y_i(t)$ moves, every other switch keeps it position:
          \item<3-> For $y_i(t)$ we have:
            \[
              x_{i-1}<f_1<f_3^-<y_i(t)\leq x_i.
            \]
            \only<3->{%
              \begin{center}
                \scriptsize
                \begin{shiftchain}[
                  column sep=3mm, row sep=3.5mm, dot radius=2.5pt,
                  labels=false, label font=\tiny,
                  bars={2/dashed,7,10/dashed,12/dashed,14},
                  bar labels={2/$x_{i-1}$,7/$y_i(t+1)$,10/\raisebox{-3.5mm}{$y_i(t)$},12/$x_i$},
                ]{16}
                  \shiftedge{3/cfred,5/cfred,8/cfnavy,16/cfred}[$\vec f$]
                \end{shiftchain}
                \par\smallskip
                $x_{i-1}<y_i(t+1)<y_i(t)\leq x_i$.
              \end{center}
            }
        \end{itemize}
      \end{proof}
    \end{column}
    \begin{column}{0.36\textwidth}
      \centering\scriptsize\setstretch{1}
      {\let\medskip\smallskip \Yexampletwo}
    \end{column}
  \end{columns}
\end{frame}

\begin{frame}[t]{Y3: the next switch stays past $e_2$}
  \footnotesize\setstretch{1}
  \tcbset{before skip=4pt,after skip=4pt}
  \claimreminder{\Xtwo
    \textbf{Auxiliary claim.} $x_{i-1}<f_1$.}
  \begin{lemma}{Claim}\Ythree\end{lemma}
  \begin{columns}[T,onlytextwidth]
    \begin{column}{0.61\textwidth}
      \begin{proof}{Proof}
        \begin{itemize}[itemsep=5pt]
          \item<1-> \textbf{Base.} If $e_1<x_j<x_{j+1}\leq e_2$,
            then $\vec e$ contains both original switches.
            \only<1->{%
              \par\smallskip
              {\centering\scriptsize
                \begin{shiftchain}[
                  column sep=4.5mm, row sep=3.5mm, dot radius=2.5pt,
                  labels=false, label font=\tiny,
                  bars={2,4,9/densely dotted},
                  bar labels={2/$x_j$,4/$x_{j+1}$,9/$e_4^-$},
                ]{10}
                \shiftedge{1/cfred,6/cfred,8/cfred,10/cfred}[$\vec e$]
                \end{shiftchain}
                \par\smallskip
                X2 requires $x_{j+1}=e_4^->e_2$.
                \par
              }
            }
          \item<2-> \textbf{Step.} Only $j=i-1$ changes.
            If $e_1<x_{i-1}<f_1$, comparability gives $\vec e<\vec f$.
            \only<2->{%
              \par\smallskip
              {\centering\scriptsize
                \begin{shiftchain}[
                  column sep=3.8mm, row sep=3.5mm, dot radius=2.5pt,
                  labels=false, label font=\tiny,
                  bars={2/dashed,7,9/dashed,11},
                  bar labels={2/$x_{i-1}$,7/$y_i(t+1)$,9/\raisebox{-3.5mm}{$y_i(t)$}},
                ]{12}
                \shiftedge{3/cfred,5/cfred,8/cfnavy,12/cfred}[$\vec f$]
                \shiftedge{1/cfred,4/cfred,7/cfred,11/cfnavy}[$\vec e$]
                \end{shiftchain}
                \par\smallskip
                $e_2\leq f_2<f_3^-=y_i(t+1)$.
                \par
              }
            }
        \end{itemize}
      \end{proof}
    \end{column}
    \begin{column}{0.36\textwidth}
      \centering\scriptsize\setstretch{1}
      {\let\medskip\smallskip \Yexamplethree}
    \end{column}
  \end{columns}
\end{frame}

\begin{frame}[t]{Y4: no $(1,2)$-edge appears}
  \footnotesize\setstretch{1}
  \tcbset{before skip=4pt,after skip=4pt}
  \claimreminder{\Xthree\Yone}
  \begin{lemma}{Claim}\Yfour\end{lemma}
  \begin{columns}[T,onlytextwidth]
    \begin{column}{0.61\textwidth}
      \begin{proof}{Proof}
        \begin{itemize}[itemsep=5pt]
          \item<1-> \textbf{Base.} X3 excludes $(1,2)$-edges in $c_{Y,0}$.
          \item<2-> \textbf{Step.} Suppose $e_1<y_j(t+1)<y_{j+1}(t+1)<e_2$.
          \item<3-> Case when none of the switches is the moved switch trivial.
          \item<4-> If $j=i-1$,
            \vspace{-1em}
            \[
              e_1<y_{i-1}(t)<f_1<f_2<y_i(t+1)=f_3^-<e_2.
            \]
            So $e_1<f_1$ but $f_2<e_2$: incomparable.
          \item<5-> If $j=i$,
            \vspace{-1em}
            \[
              e_1<y_i(t+1)<y_i(t)<y_{i+1}(t)<e_2.
            \]
            By Y1, a $(1,2)$-edge already at time $t$:
            contrary to induction.
        \end{itemize}
      \end{proof}
    \end{column}
    \begin{column}{0.36\textwidth}
      \centering\scriptsize\setstretch{1}
      % Reveal the two labelled cases alongside their proof steps.
      \begin{overlayarea}{\linewidth}{4.2cm}
        \centering
        \uncover<4->{%
          \textbf{Case $j=i-1$}\par\medskip
          \begin{shiftchain}[
            column sep=3.4mm, row sep=3.5mm, dot radius=2.5pt,
            labels=false, label font=\tiny,
            bars={3,7,9/dashed,13},
            bar labels={3/$y_{i-1}(t)$,7/$y_i(t+1)$,9/\raisebox{-3.5mm}{$y_i(t)$},13/$y_{i+1}(t)$},
          ]{14}
          \shiftedge{4/cfnavy,6/cfnavy,8/cfred,14/cfnavy}[$\vec f$]
          \shiftedge{2/cfred,10/cfred,11/cfred,12/cfred}[$\vec e$]
          \end{shiftchain}
          \par\smallskip
          $e_1<f_1$, but $f_2<e_2$.
        }
        \par\bigskip
        \uncover<5->{%
          \textbf{Case $j=i$}\par\medskip
          \begin{shiftchain}[
            column sep=3.4mm, row sep=3.5mm, dot radius=2.5pt,
            labels=false, label font=\tiny,
            bars={4,7/dashed,9},
            bar labels={4/$y_i(t+1)$,7/\raisebox{-3.5mm}{$y_i(t)$},9/$y_{i+1}(t)$},
          ]{14}
          \shiftedge{2/cfred,10/cfred,12/cfred,14/cfred}[$\vec e$]
          \end{shiftchain}
          \par\smallskip
          Both old switches were in $(e_1,e_2)$.
        }
      \end{overlayarea}
    \end{column}
  \end{columns}
\end{frame}

\begin{frame}[t]{Y5: unchanged switches keep their witness}
  \footnotesize\setstretch{1}
  \tcbset{before skip=4pt,after skip=4pt}
  \claimreminder{\Xthree\Yone\Yfive}
  \begin{lemma}{Claim}
    Y5 holds at $t=0$ and is preserved when neither switch of $\vec e$ moves.
  \end{lemma}
  \begin{proof}{Proof}
    \begin{itemize}[itemsep=4pt]
      \item<1-> \textbf{Base.} X3: no $(2,3)$-edges in $c_{Y,0}$.
      \item<2-> After the move, $\vec f$ is not monochromatic (only $f_3$ changes colour), so $\vec e\ne\vec f$.
      \item<3-> Let $\vec e$ be a $(2,3)$-edge at $t+1$ with both switches fixed.
        By Y1, it was already a $(2,3)$-edge at $t$.
      \item<4-> The next switch ($y_{j+1}(t)$ may move left, the same witness still works.
    \end{itemize}
  \end{proof}
  \begin{center}
    \only<1-3>{%
      \begin{shiftchain}[
        column sep=6mm, row sep=6mm, dot radius=3pt,
        labels=false, label font=\scriptsize,
        bars={4,7,12},
        bar labels={4/$y_{j-1}(t)$,7/$y_j(t)$,12/\raisebox{-3.5mm}{$y_{j+1}(t)$}},
      ]{16}
        \node at (8,0) {$e_3=g_3$};
        \shiftedge{1/cfred,3/cfred,8/cfred,10/cfred}[$\vec e$]
        \shiftedge{5/cfnavy,6/cfnavy,8/cfred,16/cfnavy}[$\vec g$]
      \end{shiftchain}
    }
    \only<4->{%
      \begin{shiftchain}[
        column sep=6mm, row sep=6mm, dot radius=3pt,
        labels=false, label font=\scriptsize,
        bars={4,7,10,12/dashed},
        bar labels={4/$y_{j-1}(t)$,7/$y_j(t)$,10/$y_{j+1}(t+1)$,12/\raisebox{-3.5mm}{$y_{j+1}(t)$}},
      ]{16}
        \node at (8,0) {$e_3=g_3$};
        \shiftedge{1/cfred,3/cfred,8/cfred,10/cfred}[$\vec e$]
        \shiftedge{5/cfnavy,6/cfnavy,8/cfred,16/cfnavy}[$\vec g$]
      \end{shiftchain}
    }
  \end{center}
\end{frame}

\begin{frame}[t]{Y5: moving the right switch gives witness $\vec f$}
  \footnotesize\setstretch{1}
  \tcbset{before skip=4pt,after skip=4pt}
  \claimreminder{\Yone\Ytwo\Yfive}
  \begin{lemma}{Claim}
    If the $(2,3)$-edge $\vec e$ has switches $y_{i-1}(t),y_i(t+1)$,
    then $e_3=f_3$ and $\vec g=\vec f$ is a witness.
  \end{lemma}
  \begin{proof}{Proof}
    \begin{itemize}[itemsep=4pt]
      \item<1-> Y1 puts the next switch outside $\vec e$:
        $e_4<y_{i+1}(t)<f_4$.
      \item<2-> Comparability gives $e_3\leq f_3$;
        $f_3^-=y_i(t+1)<e_3$ gives $f_3\leq e_3$. Thus $e_3=f_3$.
      \item<3-> Take $\vec g=\vec f$:
        $e_3^-=y_i(t+1)=g_3^-<y_i(t)\leq x_i$,
        and $y_{i+1}(t+1)=y_{i+1}(t)<g_4$.
    \end{itemize}
  \end{proof}
  \begin{center}
    \begin{shiftchain}[
      column sep=6mm, row sep=6mm, dot radius=3pt,
      labels=false, label font=\scriptsize,
      bars={4,8,10/dashed,13},
      bar labels={4/$y_{i-1}(t)$,8/$y_i(t+1)$,10/\raisebox{-3.5mm}{$y_i(t)$},13/$y_{i+1}(t)$},
    ]{16}
      \uncover<2->{\node at (9,0) {$e_3=f_3$};}
      \shiftedge{1/cfred,3/cfred,9/cfred,12/cfred}[$\vec e$]
      \shiftedge{5/cfnavy,7/cfnavy,9/cfred,16/cfnavy}[$\vec f$]
    \end{shiftchain}
  \end{center}
\end{frame}

\begin{frame}[t]{Y5: moving the left switch keeps the old witness}
  \footnotesize\setstretch{1}
  \tcbset{before skip=4pt,after skip=4pt}
  \claimreminder{\Yone\Yfive}
  \begin{lemma}{Claim}
    If the $(2,3)$-edge $\vec e$ has switches $y_i(t+1),y_{i+1}(t)$,
    its old witness still works.
  \end{lemma}
  \begin{proof}{Proof}
    \begin{itemize}[itemsep=4pt]
      \item<1-> Both old switches were already in the middle gap:
        $e_2<y_i(t+1)<y_i(t)<y_{i+1}(t)<e_3$.
      \item<2-> Y1 excludes any further switch at $t$;
        $\vec e$ was a $(2,3)$-edge, so induction gives $\vec g$.
      \item<3-> The right switch and the next switch stay fixed:
        $e_3^-=y_{i+1}(t+1)=g_3^-<x_{i+1}$,
        and $y_{i+2}(t+1)=y_{i+2}(t)<g_4$.
    \end{itemize}
  \end{proof}
  \begin{center}
    \begin{shiftchain}[
      column sep=6mm, row sep=6mm, dot radius=3pt,
      labels=false, label font=\scriptsize,
      bars={4,6/dashed,9,13},
      bar labels={4/$y_i(t+1)$,6/\raisebox{-3.5mm}{$y_i(t)$},9/$y_{i+1}(t)$,13/$y_{i+2}(t)$},
    ]{16}
      \node at (10,0) {$e_3=g_3$};
      \shiftedge{1/cfred,3/cfred,10/cfred,12/cfred}[$\vec e$]
      \shiftedge{7/cfnavy,8/cfnavy,10/cfred,16/cfnavy}[$\vec g$]
    \end{shiftchain}
  \end{center}
\end{frame}

\begin{frame}[t]{Y6: chosen $(3,4)$-edges move left}
  \small\setstretch{1.05}
  \claimreminder{\Yone}
  \begin{lemma}{Claim}$\vec f(t+1)<\vec f(t)$.\end{lemma}
  \begin{proof}{Proof}
    \begin{itemize}
      \item<1-> Suppose the next chosen edge $\vec e>\vec f$. It must be new, since $\vec f$ was maximal.
      \item<2-> $y_i(t+1)=f_3^-<f_3\leq e_3$. A $(3,4)$-edge has both switches after $e_3$, so the new bar is outside $\vec e$.
      \item<3-> If the old bar was outside too, $\vec e$ was already monochromatic. If inside, Y1 leaves at most one switch. Both are impossible.
    \end{itemize}
\only<2->{\begin{center}
\begin{shiftchain}[column sep=5mm,row sep=4mm,dot radius=3.0pt,labels=false,label font=\scriptsize,
bars={5,8/dashed,11,14},bar labels={5/$f_3^-$,8/$y_i(t)$,11/$y_{i+1}$,14/$y_{i+2}$}]{16}
\shiftedge{1/cfred,3/cfred,6/cfnavy,13/cfred}[$\vec f$]
\shiftedge{7/cfnavy,9/cfnavy,10/cfnavy,14/cfred}[$\vec e\text{ (hypothetical)}$]
\end{shiftchain}
\par\smallskip{\scriptsize The moved bar lies left of $e_1$: it cannot supply a new pair in $(e_3,e_4)$.}
\end{center}}
  \end{proof}
\end{frame}

\begin{frame}[t]{Y6: termination leaves only $(2,3)$-edges}
  \small\setstretch{1.05}
  \claimreminder{\Yone\Yfour\textbf{Previous step.} $\vec f(t+1)<\vec f(t)$.}
  \begin{lemma}{Claim}\Ysix\end{lemma}
  \begin{proof}{Proof}
    \begin{itemize}
      \item<1-> By the previous step, the chosen edges strictly decrease, so Stage Y terminates.
      \item<2-> At termination: no $(3,4)$-edge; Y4 excludes $(1,2)$-edges. By Y1, only $(2,3)$-edges can be monochromatic.
    \end{itemize}
\only<2->{\begin{center}
\begin{shiftchain}[column sep=5mm,row sep=4mm,dot radius=3.0pt,labels=false,label font=\scriptsize,
bars={5,8},bar labels={5/$y_{j-1}$,8/$y_j$}]{16}
\shiftedge{1/cfred,3/cfred,10/cfred,13/cfred}[$(2,3)$]
\end{shiftchain}
\par\smallskip{\scriptsize The only surviving monochromatic pattern: two bars in the middle gap.}
\end{center}}
  \end{proof}
\end{frame}

